% from ega1-10-fr.tex
\oldpage[I]{206}

\subsection{Morphismes adiques de préschémas formels.}
\label{subsection:I.10.12-fr}

\begin{env}[10.12.1]
\label{I.10.12.1-fr}
Soient $\mathfrak{X}$, $\mathfrak{S}$ deux préschémas formels
\emph{localement noethériens}~; nous dirons qu'un morphisme
$f:\mathfrak{X}\to\mathfrak{S}$ est \emph{adique} s'il existe un Idéal
de définition $\mathscr{J}$ de $\mathfrak{S}$ tel que
$\mathscr{K}=f^*(\mathscr{J})\mathscr{O}_{\mathfrak{X}}$ soit un Idéal de
définition de $\mathfrak{X}$~; on dit aussi alors que $\mathfrak{X}$ est
un \emph{$\mathfrak{S}$-préschéma adique} (pour $f$). Lorsqu'il en est
ainsi, pour \emph{tout} Idéal de définition $\mathscr{J}_1$ de
$\mathfrak{S}$,
$\mathscr{K}_1=f^*(\mathscr{J}_1)\mathscr{O}_{\mathfrak{X}}$ est un Idéal
de définition de $\mathfrak{X}$. En effet, la question étant locale, on
peut supposer $\mathfrak{X}$ et $\mathfrak{S}$ affines noethériens~; il
existe donc un entier $n$ tel que $\mathscr{J}^n\subset\mathscr{J}_1$ et
$\mathscr{J}_1^n\subset\mathscr{J}$
(\hyperref[I.10.3.6-fr]{10.3.6} et
\hyperref[0.7.1.4-fr]{0, 7.1.4}), d'où
$\mathscr{K}^n\subset\mathscr{K}_1$ et
$\mathscr{K}_1^n\subset\mathscr{K}$. La première de ces relations prouve
que $\mathscr{K}_1=\mathfrak{K}_1^\Delta$, où $\mathfrak{K}_1$ est un
idéal ouvert de
$A=\Gamma(\mathfrak{X},\mathscr{O}_{\mathfrak{X}})$, et la seconde prouve
que $\mathfrak{K}_1$ est un idéal de définition de $A$
(\hyperref[0.7.1.4-fr]{0, 7.1.4}), d'où notre assertion.
\end{env}

Il résulte aussitôt de ce qui précède que si $\mathfrak{X}$ et
$\mathfrak{Y}$ sont deux $\mathfrak{S}$-préschémas adiques, \emph{tout
$\mathfrak{S}$-morphisme $u:\mathfrak{X}\to\mathfrak{Y}$ est adique}~:
en effet, si $f:\mathfrak{X}\to\mathfrak{S}$,
$g:\mathfrak{Y}\to\mathfrak{S}$ sont les morphismes structuraux, et
$\mathscr{J}$ un Idéal de définition de $\mathfrak{S}$, on a
$f=g\circ u$, donc
$u^*(g^*(\mathscr{J})\mathscr{O}_{\mathfrak{Y}})
\mathscr{O}_{\mathfrak{X}}=f^*(\mathscr{J})
\mathscr{O}_{\mathfrak{X}}$ est un Idéal de définition de
$\mathfrak{X}$, et par hypothèse
$g^*(\mathscr{J})\mathscr{O}_{\mathfrak{Y}}$ est un Idéal de définition
de $\mathfrak{Y}$.

\begin{env}[10.12.2]
\label{I.10.12.2-fr}
Dans ce qui suit, nous supposerons fixé un préschéma formel localement
noethérien $\mathfrak{S}$ et un Idéal de définition $\mathscr{J}$ de
$\mathfrak{S}$~; nous poserons
$S_n=(\mathfrak{S},\mathscr{O}_{\mathfrak{S}}/\mathscr{J}^{n+1})$. Les
$\mathfrak{S}$-préschémas adiques (localement noethériens) forment
évidemment une \emph{catégorie}. Nous dirons qu'un système inductif
$(X_n)$ de $S_n$-préschémas (usuels) localement noethériens est un
\emph{$(S_n)$-système inductif adique} si les morphismes structuraux
$f_n:X_n\to S_n$ sont tels que, pour $m\leq n$, les diagrammes
\[
  \label{I.10.12.2.1-fr}
  \xymatrix{
    X_n \ar[d]_{f_n}
    & X_m \ar[l] \ar[d]^{f_m}\\
    S_n
    & S_m \ar[l]
  }
  \tag{10.12.2.1}
\]
soient commutatifs et \emph{identifient $X_m$ au produit
$X_n\times_{S_n}S_m=(X_n)_{(S_m)}$}. Les systèmes inductifs adiques
forment une \emph{catégorie}~: il suffit en effet de définir un morphisme
$(X_n)\to(Y_n)$ de tels systèmes comme un \emph{système inductif de
$S_n$-morphismes $u_n:X_n\to Y_n$} tel que $u_m$ s'identifie à
$(u_n)_{(S_m)}$ pour $m\leq n$. Cela étant~:
\end{env}

\begin{theorem}[10.12.3]
\label{I.10.12.3-fr}
Il y a équivalence canonique entre la catégorie des
$\mathfrak{S}$-préschémas adiques et la catégorie des
$(S_n)$-systèmes inductifs adiques.
\end{theorem}

L'équivalence en question s'obtient de la façon suivante~: si
$\mathfrak{X}$ est un $\mathfrak{S}$-préschéma adique,
$f:\mathfrak{X}\to\mathfrak{S}$ le morphisme structural,
$\mathscr{K}=f^*(\mathscr{J})\mathscr{O}_{\mathfrak{X}}$ est un Idéal de
définition de $\mathfrak{X}$ et on fait correspondre à $\mathfrak{X}$ le
système inductif des
$X_n=(\mathfrak{X},\mathscr{O}_{\mathfrak{X}}/\mathscr{K}^{n+1})$, le
morphisme structural $f_n:X_n\to S_n$ correspondant à $f$
(\hyperref[I.10.5.6-fr]{10.5.6}). Montrons d'abord que $(X_n)$ est un
\emph{système inductif adique}~: si $f=(\psi,\theta)$, on a
$\psi^*(\mathscr{J})\mathscr{O}_{\mathfrak{X}}=\mathscr{K}$, donc
$\psi^*(\mathscr{J}^n)\mathscr{O}_{\mathfrak{X}}=\mathscr{K}^n$ pour tout
$n$, et (par l'exactitude du foncteur $\psi^*$)
$\mathscr{K}^{m+1}/\mathscr{K}^{n+1}
=\psi^*(\mathscr{J}^{m+1}/\mathscr{J}^{n+1})
(\mathscr{O}_{\mathfrak{X}}/\mathscr{K}^{n+1})$ pour $m\leq n$~; notre
conclusion résulte donc de (\hyperref[I.4.4.5-fr]{4.4.5}). Il est immédiat
en outre de vérifier qu'à un $\mathfrak{S}$-morphisme
$u:\mathfrak{X}\to\mathfrak{Y}$ de $\mathfrak{S}$-préschémas adiques
correspond (avec des notations évidentes)

